\documentclass[a4paper,12pt]{article}
\usepackage{ctex}
\usepackage{geometry}
\usepackage{listings}
\usepackage{xcolor}
\usepackage{enumitem}

\geometry{left=2.5cm, right=2.5cm, top=2.5cm, bottom=2.5cm}

\definecolor{codegreen}{rgb}{0,0.6,0}
\definecolor{codegray}{rgb}{0.5,0.5,0.5}
\definecolor{codepurple}{rgb}{0.58,0,0.82}
\definecolor{backcolour}{rgb}{0.95,0.95,0.92}

\lstdefinestyle{mystyle}{
    backgroundcolor=\color{backcolour},   
    commentstyle=\color{codegreen},
    keywordstyle=\color{magenta},
    numberstyle=\tiny\color{codegray},
    stringstyle=\color{codepurple},
    basicstyle=\ttfamily\footnotesize,
    breakatwhitespace=false,         
    breaklines=true,                 
    captionpos=b,                    
    keepspaces=true,                 
    numbers=left,                    
    numbersep=5pt,                  
    showspaces=false,                
    showstringspaces=false,
    showtabs=false,                  
    tabsize=2
}

\lstset{style=mystyle}

\title{专升本 C 语言题库整理}
\author{}
\date{}

\begin{document}

\maketitle

\section{文件一：专升本 c 语言题库及答案}

\subsection{一、选择题（每题 2 分，共 40 分）}

\begin{enumerate}
    \item C 语言中，以下哪个选项是合法的数据类型？\\
    A) int \\
    B) float \\
    C) char \\
    D) All of the above \\
    答案：D

    \item 在 C 语言中，以下哪个符号表示逻辑与运算符？\\
    B) || \\
    C) ! \\
    D) | \\
    答案：A

    \item 如果一个整型变量声明为 int a = 5;，那么 a 的值是多少？\\
    A) 5 \\
    B) 0 \\
    C) 1 \\
    D) 无法确定 \\
    答案：A

    \item 在 C 语言中，以下哪个选项是正确的注释方式？\\
    A) // \\
    C) // \\
    答案：B

    \item 下列哪个函数用于动态分配内存？\\
    A) malloc() \\
    B) alloc() \\
    C) new() \\
    D) calloc() \\
    答案：A

    \item 在 C 语言中，以下哪个选项表示字符串结束符？\\
    A) '\textbackslash 0' \\
    B) '\textbackslash n' \\
    C) '\textbackslash r' \\
    D) '\textbackslash t' \\
    答案：A

    \item 在 C 语言中，以下哪个函数用于输入字符串？\\
    A) scanf() \\
    B) printf() \\
    C) gets() \\
    D) fgets() \\
    答案：C

    \item 在 C 语言中，以下哪个函数用于输出字符串？\\
    A) scanf() \\
    B) printf() \\
    C) gets() \\
    D) fgets() \\
    答案：B

    \item 在 C 语言中，以下哪个选项表示结构体类型？\\
    A) struct \\
    B) union \\
    C) enum \\
    D) class \\
    答案：A

    \item 在 C 语言中，以下哪个选项表示共用体类型？\\
    A) struct \\
    B) union \\
    C) enum \\
    D) class \\
    答案：B

    \item 在 C 语言中，以下哪个选项表示枚举类型？\\
    A) struct \\
    B) union \\
    C) enum \\
    D) class \\
    答案：C

    \item 在 C 语言中，以下哪个选项表示类类型？\\
    A) struct \\
    B) union \\
    C) enum \\
    D) class \\
    答案：D

    \item 下列哪个函数用于打开文件？\\
    A) fopen() \\
    B) fclose() \\
    C) fread() \\
    D) fwrite() \\
    答案：A

    \item 下列哪个函数用于关闭文件？\\
    A) fopen() \\
    B) fclose() \\
    C) fread() \\
    D) fwrite() \\
    答案：B

    \item 在 C 语言中，以下哪个选项表示函数返回类型为 void？\\
    A) int \\
    B) float \\
    C) char \\
    D) void \\
    答案：D

    \item 在 C 语言中，以下哪个选项表示函数参数列表为空？\\
    A) int \\
    B) void \\
    C) char \\
    D) () \\
    答案：D

    \item 在 C 语言中，以下哪个选项表示函数的返回值类型为整型？\\
    A) int \\
    B) float \\
    C) char \\
    D) void \\
    答案：A

    \item 在 C 语言中，以下哪个选项表示函数的返回值类型为浮点型？\\
    A) int \\
    B) float \\
    C) char \\
    D) void \\
    答案：B

    \item 在 C 语言中，以下哪个选项表示函数的返回值类型为字符型？\\
    A) int \\
    B) float \\
    C) char \\
    D) void \\
    答案：C

    \item 在 C 语言中，以下哪个选项表示函数的返回值类型为 void？\\
    A) int \\
    B) float \\
    C) char \\
    D) void \\
    答案：D
\end{enumerate}

\subsection{二、填空题（每题 2 分，共 40 分）}

\begin{enumerate}
    \item 在 C 语言中，一个整型变量占用的内存大小为\_\_\_\_\_\_。\\
    答案：4 字节

    \item 在 C 语言中，一个字符型变量占用的内存大小为\_\_\_\_\_\_。\\
    答案：1 字节

    \item 在 C 语言中，一个浮点型变量占用的内存大小为\_\_\_\_\_\_。\\
    答案：4 字节

    \item 在 C 语言中，一个双精度浮点型变量占用的内存大小为\_\_\_\_\_\_。\\
    答案：8 字节

    \item 在 C 语言中，一个 short 类型变量占用的内存大小为\_\_\_\_\_\_。\\
    答案：2 字节

    \item 在 C 语言中，一个 long 类型变量占用的内存大小为\_\_\_\_\_\_。\\
    答案：4 字节或 8 字节（取决于操作系统）

    \item 在 C 语言中，一个 unsigned char 类型变量占用的内存大小为\_\_\_\_\_\_。\\
    答案：1 字节

    \item 在 C 语言中，一个 unsigned int 类型变量占用的内存大小为\_\_\_\_\_\_。\\
    答案：4 字节

    \item 在 C 语言中，一个 unsigned long 类型变量占用的内存大小为\_\_\_\_\_\_。\\
    答案：4 字节或 8 字节（取决于操作系统）

    \item 在 C 语言中，一个指针变量占用的内存大小为\_\_\_\_\_\_。\\
    答案：4 字节或 8 字节（取决于操作系统）

    \item 在 C 语言中，以下程序段的输出结果是\_\_\_\_\_\_。
    \begin{lstlisting}[language=C]
int main() {
    int a = 5;
    printf("%d\n", ++a);
    return 0;
}
    \end{lstlisting}
    答案：6

    \item 在 C 语言中，以下程序段的输出结果是\_\_\_\_\_\_。
    \begin{lstlisting}[language=C]
int main() {
    int a = 5;
    printf("%d\n", a++);
    return 0;
}
    \end{lstlisting}
    答案：5

    \item 在 C 语言中，以下程序段的输出结果是\_\_\_\_\_\_。
    \begin{lstlisting}[language=C]
int main() {
    int a = 5;
    printf("%d\n", a + 5);
    return 0;
}
    \end{lstlisting}
    答案：10

    \item 在 C 语言中，以下程序段的输出结果是\_\_\_\_\_\_。
    \begin{lstlisting}[language=C]
int main() {
    int a = 5;
    printf("%d\n", a / 2);
    return 0;
}
    \end{lstlisting}
    答案：2

    \item 在 C 语言中，以下程序段的输出结果是\_\_\_\_\_\_。
    \begin{lstlisting}[language=C]
int main() {
    int a = 5;
    printf("%d\n", a \% 2);
    return 0;
}
    \end{lstlisting}
    答案：1

    \item 在 C 语言中，以下程序段的输出结果是\_\_\_\_\_\_。
    \begin{lstlisting}[language=C]
int main() {
    int a = 5;
    printf("%d\n", a > 4);
    return 0;
}
    \end{lstlisting}
    答案：1

    \item 在 C 语言中，以下程序段的输出结果是\_\_\_\_\_\_。
    \begin{lstlisting}[language=C]
int main() {
    int a = 5;
    printf("%d\n", a < 6);
    return 0;
}
    \end{lstlisting}
    答案：1

    \item 在 C 语言中，以下程序段的输出结果是\_\_\_\_\_\_。
    \begin{lstlisting}[language=C]
int main() {
    int a = 5;
    printf("%d\n", a >= 5);
    return 0;
}
    \end{lstlisting}
    答案：1

    \item 在 C 语言中，以下程序段的输出结果是\_\_\_\_\_\_。
    \begin{lstlisting}[language=C]
int main() {
    int a = 5;
    printf("%d\n", a <= 5);
    return 0;
}
    \end{lstlisting}
    答案：1

    \item 在 C 语言中，以下程序段的输出结果是\_\_\_\_\_\_。
    \begin{lstlisting}[language=C]
int main() {
    int a = 5;
    printf("%d\n", a == 5);
    return 0;
}
    \end{lstlisting}
    答案：1
\end{enumerate}

\subsection{三、编程题（每题 20 分，共 60 分）}

\begin{enumerate}
    \item 编写一个 C 程序，实现以下功能：从键盘输入 10 个整数，计算并输出这 10 个整数的平均值。\\
    答案：
    \begin{lstlisting}[language=C]
int main() {
    int numbers[10];
    int sum = 0;
    float average;
    printf("请输入 10 个整数：\n");
    for (int i = 0; i < 10; i++) {
        sum += numbers[i];
    }
    average = sum / 10.0;
    printf("这 10 个整数的平均值为：%f\n", average);
    return 0;
}
    \end{lstlisting}

    \item 编写一个 C 程序，实现以下功能：从键盘输入一个字符串，计算并输出这个字符串的长度。\\
    答案：
    \begin{lstlisting}[language=C]
int main() {
    char str[100];
    printf("请输入一个字符串：\n");
    gets(str);
    int length = strlen(str);
    printf("这个字符串的长度为：%d\n", length);
    return 0;
}
    \end{lstlisting}

    \item 编写一个 C 程序，实现以下功能：从键盘输入一个整数，判断这个整数是否为素数，并输出结果。\\
    答案：
    \begin{lstlisting}[language=C]
bool is_prime(int num) {
    if (num <= 1) {
        return false;
    }
        if (num % i == 0) {
            return false;
        }
    }
    return true;
}
int main() {
    int num;
    printf("请输入一个整数：\n");
    if (is_prime(num)) {
        printf("%d 是素数\n", num);
    } else {
        printf("%d 不是素数\n", num);
    }
    return 0;
}
    \end{lstlisting}
\end{enumerate}

\section{文件二：C 专升本复习 - 考试题型及基本算法罗列}

\subsection{考试题型}
\begin{itemize}
    \item 一、选择（20\% 共 20 题）
    \item 二、填空（20\% 共 20 题
    \item 三、程序填空（10\% 共 5 题）
    \item 四、看程序写结果（20\% 共 5 题）
    \item 五、改错（10\% 共 2 题）
    \item 六、编程（20\% 共 2 题）
\end{itemize}

\subsection{基本算法罗列}

\begin{enumerate}
    \item 打印出 100 到 200 之间的所有素数（即除 1 和它本身再没有其他约数的数）。
    \begin{lstlisting}[language=C++]
#include <iostream.h>
#include <math.h>
void main()
{
int i,j;
for (i=100;i<=200;i++)
{
int temp=int(sqrt(i));
for (j=2;      (1)      ;j++)
if (i%j==0)          (2)        ;
if (      (3)      )   cout<<i<<' ';
}
cout<<'\n';
}
(1)    j<=temp          (2)    break        (3) j>temp
    \end{lstlisting}

    \item 斐波那契数列的第 1 和第 2 个数分别为 0 和 1 ，从第三个数开始，每个数等于其前两个数之和。求斐波那契数列中的前 20 个数，要求每行输出 5 个数。
    \begin{lstlisting}[language=C++]
#include<iostream.h>
void main() {
int f,f1,f2,i;
cout<<" 斐波那契数列：\n";
f1=0; f2=1;
cout<<setw(6)<<f1<<setw(6)<<f2;
for(i=3;i<=20;i++) {
f=______(1)______; 
cout<<setw(6)<<f;
if(_____(2)______) cout<<endl;
f1=f2;
f2=____(3)_______;
}
cout<<endl;
}
(1)  f1+f2      (2) i%5==0      (3) f
    \end{lstlisting}

    \item 计算
    的值。
    \begin{lstlisting}[language=C++]
#include<iostream.h>
void main()
{
double x,p1=1,p2=1,s=0;
int i,j=1;
cout<<"输入 x 的值:";
cin>>x;
for(i=1;i<=10;i++) {
p1*=___(1)_____;
p2*=____(2)____;
s+=j*p1/p2;  //j 的值为 (-1)i+1
j=____(3)____;
}
cout<<s<<endl;
}
(1) x                     (2) I            (3) –j                     
    \end{lstlisting}

    \item 采用辗转相除法求出两个整数的最大公约数。
    \begin{lstlisting}[language=C++]
#include<iostream.h>
void main()
{
int a,b;
cout<<"请输入两个正整数:";
cin>>a>>b;
while(a<=0 || __(1)___) {cout<<"重新输入:"; cin>>a>>b;}
while(b) {
int r;
r=a%b;
___(2)___; ___(3)___;  //分别修改 a 和 b 的值
}
cout<<a<<endl;  //输出最大公约数
}
(1) b<=0                     (2) a=b          (3) b=r
    \end{lstlisting}

    \item 把从键盘上输入的一个大于等于 3 的整数分解为质因子的乘积。如输入 24 时得到的输出结果为“2 2 2 3”，输入 50 时得到的输出结果为“2 5 5”，输入 37 时得到的输出结果为“37”。
    \begin{lstlisting}[language=C++]
#include<iostream.h>
void main()
{
int x;
cout<<"请输入一个整数，若小于 3 则重输:";
do cin>>x; while(___(1)___);
int i=2;
do{
while(___(2)___) {
cout<<i<<' ';
x/=i;
}
___(3)___;
}while(i<x);
if(x!=1) cout<<x; 
cout<<endl;
}
(1) x<3 (或 x<=2)            (2) x%i==0      (3) i++
    \end{lstlisting}

    \item 下面函数是求两个整型参数 a 和 b 的最小公倍数。
    \begin{lstlisting}[language=C++]
int f2(int a, int b)
{
int i=2, p=1;
do {
while(a%i==0 && ___(1)___) {
p*=i; a/=i; b/=i;
}
___(2)___;
}while(a>=i && ___(3)___);
return p*a*b;
}
(1) b%i==0                   (2) i++ (或++i)  (3) b>=i
    \end{lstlisting}

    \item 在输出屏幕上打印出一个由字符'*'组成的等腰三角形，该三角形的高为 5 行，从上到下每行的字符数依次为 1,3,5,7,9。
    \begin{lstlisting}[language=C++]
#include<iostream.h>
void main()
{
int i,j;
for(i=1;___(1)___;i++) {
for(j=1;j<=9;j++)
if(j<=5-i || ___(2)___) cout<<' ';
else ___(3)___;
cout<<endl;
}
}
(1) i<=5                     (2) j>=5+i        (3) cout<<'*'
    \end{lstlisting}

    \item 统计字符串中英文字母个数的程序。
    \begin{lstlisting}[language=C++]
#include <iostream.h>
int count (char str[]);
void main(){
char s1[80];
cout <<"Enter a line:";
cin >>s1;
cout <<"count="<<count(s1)<<endl;
}
int count(char str[]){
int num=0;  //给统计变量赋初值
for(int i=0;str[i];i++)
if (str[i]>='a' && str[i]<='z' ||___(1)___ )
___(2)___;
___(3)___;
}
(1) str[i]>='A' && str[i]<='Z'  (2) num++    (3) return num 
    \end{lstlisting}

    \item 主函数调用一个 fun 函数将字符串逆序。
    \begin{lstlisting}[language=C++]
#include<iostream.h>
#include<string.h>
___(1)___;
void main( ) {
char s[80];
cin>>s;
___(2)___;
cout<<"逆序后的字符串:"<<s<<endl ;
}
void fun(char ss[]) {
int n=strlen(ss);
for(int i=0; ___(3)____; i++) {
char c=ss[i];
ss[i]=ss[n–1–i];
ss[n–1–i]=c;
}
}
(1) void fun(char ss[])    (2) fun(s)       (3) i<n/2
    \end{lstlisting}

    \item 从一个字符串中删除所有同一个给定字符后得到一个新字符串并输出。
    \begin{lstlisting}[language=C++]
#include<iostream.h>
const int len=20;
void delstr(char a[],char b[],char c);
void main() {
char str1[len],str2[len];
char ch;
cout<<"输入一个字符串:";
cin>>str1;
cout<<"输入一个待删除的字符:"; 
cin>>ch;
delstr(str1,str2,ch);
cout<<str2<<endl;
}
void delstr(char a[],char b[],char c)
{
int j=0;
for(int i=0; ___(1)___; i++)
    \end{lstlisting}
\end{enumerate}

\section{文件三：C 语言专升本 100 题}

\subsection{一、填空题}
\begin{enumerate}
    \item 在 C 语言中，一个函数一般由两个部分组成，它们分别是\_\_\_\_\_\_\_\_和\_\_\_\_\_\_\_\_。
    \item 若有定义：int a[10],*p=a;，则*(p+5) 表示\_\_\_\_\_\_\_\_。
    \item 设 x、y、z 均为 int 型变量，则执行以下语句后，z 的值为\_\_\_\_\_\_\_\_。\\
    x=y=z=1;++x||++y\&\&++z;
    \item 若有以下定义和语句：\\
    char c1='b',c2='e';\\
    printf("\%d,\%c\n",c2-c1,c2-'a'+'A');\\
    则输出结果是\_\_\_\_\_\_\_\_。
    \item 以下程序的输出结果是\_\_\_\_\_\_\_\_。
    \begin{lstlisting}[language=C]
include <stdio.h>
void main( )
{
 int a=1,b=4,c=2;
 if(a<b)
 if(b<0) c=0;
 else c++;
 printf("%d\n",c);
}
    \end{lstlisting}
\end{enumerate}

\subsection{二、选择题}
\begin{enumerate}
    \item 若有以下定义：\\
    char s[10]="hello\textbackslash 0\textbackslash t\textbackslash \textbackslash \textbackslash \textbackslash ";\\
    则 printf("\%d\n",strlen(s)); 的输出结果是（ ）。\\
    A. 10 \\
    B. 9 \\
    C. 7 \\
    D. 6

    \item 下面程序的运行结果是（ ）。
    \begin{lstlisting}[language=C]
include <stdio.h>
void main()
{
 int x=1,y=0,a=0,b=0;
 switch(x)
 {
 case 1:
 switch(y)
 {
 case 0: a++; break;
 case 1: b++; break;
 }
 case 2: a++; b++; break;
 case 3: a++; b++;
 }
 printf("a=%d,b=%d\n",a,b);
}
    \end{lstlisting}
    A. a=2,b=1 \\
    B. a=1,b=1 \\
    C. a=1,b=0 \\
    D. a=2,b=2

    \item 有以下程序
    \begin{lstlisting}[language=C]
include <stdio.h>
void main()
{
 int a=1,b=0;
 if(!a) b++;
 else if(a==0)
 if(a) b++;
 else b++;
 printf("%d\n",b);
}
    \end{lstlisting}
    程序运行后的输出结果是（ ）。\\
    A. 0 \\
    B. 1 \\
    C. 2 \\
    D. 3

    \item 以下程序的输出结果是（ ）。
    \begin{lstlisting}[language=C]
include <stdio.h>
void main()
{
 int a=5,b=4,c=3,d;
 d=(a>b>c);
 printf("%d\n",d);
}
    \end{lstlisting}
    A. 0 \\
    B. 1 \\
    C. 非 0 的数 \\
    D. -1

    \item 以下程序的输出结果是（ ）。
    \begin{lstlisting}[language=C]
include <stdio.h>
void main()
{
 int a=0,b=1,c=0,d=20;
 if(a) d=d-10;
 else if(!b)
 if(!c) d=15;
 else d=25;
 printf("%d\n",d);
}
    \end{lstlisting}
    A. 10 \\
    B. 20 \\
    C. 30 \\
    D. 程序进入死循环

    \item 若有以下定义和语句：\\
    char c1='b',c2='e';\\
    printf("\%d,\%c\n",c2-c1,c2-'a'+'A');\\
    则输出结果是（ ）。\\
    A. 2, M \\
    B. 3, E \\
    C. 2, E \\
    D. 输出项与对应的格式控制不一致，输出结果不确定

    \item 若有以下说明和语句：\\
    int c[4][5],(*p)[5];\\
    p=c;\\
    能够正确引用数组元素 c[1][2] 的是（ ）。\\
    A. *(*(p+1)+2) \\
    B. (*p)[2] \\
    C. p+1[2] \\
    D. *(p+5)

    \item 以下程序的输出结果是（ ）。
    \begin{lstlisting}[language=C]
include <stdio.h>
void main()
{
 int a[10]={1,2,3,4,5,6,7,8,9,10},*p=\&a[3],b;
 b=*(p+2);
 printf("%d\n",b);
}
    \end{lstlisting}
    A. 6 \\
    B. 7 \\
    C. 8 \\
    D. 9

    \item 若有以下定义和语句：\\
    int a[10]=\{1,2,3,4,5,6,7,8,9,10\},*p=a;\\
    则值为 6 的表达式是（ ）。\\
    A. *p+6 \\
    B. *p++ \\
    C. p+6 \\
    D. p+=5,*(p+5)

    \item 有以下程序
    \begin{lstlisting}[language=C]
include <stdio.h>
void main()
{
 int a[3][3]=\{\{1,2,3\},\{4,5,6\},\{7,8,9\}\};
 int i,j,s=0;
 for(i=0;i<3;i++)
 for(j=0;j<i;j++) s+=a[i][j];
 printf("%d\n",s);
}
    \end{lstlisting}
    程序运行后的输出结果是（ ）。\\
    A. 18 \\
    B. 19 \\
    C. 20 \\
    D. 21

    \item 以下程序的输出结果是（ ）。
    \begin{lstlisting}[language=C]
include <stdio.h>
void main()
{
 int a[3][3]=\{\{1,2,3\},\{4,5,6\},\{7,8,9\}\};
 int i,j,s=0;
 for(i=0;i<3;i++)
 for(j=0;j<=i;j++) s+=a[i][j];
 printf("%d\n",s);
}
    \end{lstlisting}
    A. 18 \\
    B. 19 \\
    C. 20 \\
    D. 21

    \item 有以下程序
    \begin{lstlisting}[language=C]
include <stdio.h>
void fun(char *c, int d)
{
 *c=*c+1; d=d+1;
 printf("%c, %c, ", *c, d);
}
void main()
{
 char a='A', b='a';
 fun(\&b, a);
 printf("%c, %c\n", a, b);
}
    \end{lstlisting}
    程序运行后的输出结果是（ ）。\\
    A. B, b, A, a \\
    B. a, B, B, b \\
    C. A, b, B, a \\
    D. b, B, a, B

    \item 若有以下程序段：\\
    int a[4]=\{0,0,0,0\},i;\\
    for(i=1;i<4;i++) scanf("\%d", \&a[i]);\\
    若从键盘输入：1234<回车>，则 a[1] 的值是（ ）。\\
    A. 不定值 \\
    B. 0 \\
    C. 1234 \\
    D. 4
\end{enumerate}

\subsection{三、简答题}
\begin{enumerate}
    \item 请简要介绍 C 语言中的数据类型。
    \item 请简要介绍 C 语言中的指针。
    \item 请简要介绍 C 语言中的数组。
\end{enumerate}

\subsection{四、编程题}
\begin{enumerate}
    \item 编写一个程序，从键盘输入一个整数，判断该整数是否为素数。
    \item 编写一个程序，计算 1 到 n 的阶乘之和。
    \item 编写一个程序，输入一个字符串，将其转换为大写字母并输出。
    \item 编写一个程序，输入一个整数，将其转换为二进制字符串并输出。
\end{enumerate}

\end{document}